\chapter*{GENERAL CONCLUSION}
\markboth{\MakeUppercase{GENERAL CONCLUSION}}{}
\addcontentsline{toc}{chapter}{GENERAL CONCLUSION}
\adjustmtc
\thispagestyle{MyStyle}

The work presented in this report set out to address a structural fragility observed during a five-day audit at \clientCompany: an estimation organisation rich in expertise and discipline at the individual level, but underserved by the connective tissue --- the visibility, the traceability, the historical learning, the early surfacing of feasibility risk --- that would let that expertise compose into a coordinated process. The findings motivated a reframing of the original mandate. Rather than pursue the automation of the Bill of Quantities, a goal the audit found infeasible under the cost-of-error conditions of industrial quotation, the project committed to build a platform that strengthens every part of the process that does not depend on irreducible human expertise, and respects the part that does. The resulting platform, Itqan, is the subject of this report.

\textit{The strategic reframing itself.} The pivot from BOQ automation to RFQ Lifecycle Intelligence (Chapter~1) is the foundational contribution: every subsequent contribution rests on that reframing.

\textit{The four-pillar target architecture.} Workflow and Tracking is substantially delivered as the operational backbone; Communication Automation and Executive/Business Intelligence are partial, addressed through reminder infrastructure, analytical endpoints, and demonstration dashboards; Historical Learning and Prediction is delivered at the foundation level through deterministic document parsing, with the predictive components left as future work.

\textit{The trust-boundary architecture of the RFQ Copilot.} The copilot architecture developed in Chapter~4 is the central engineering contribution. It allocates classification, policy enforcement, evidence selection, response composition, guardrail verification, and language-level judgment to distinct components, separates classification from execution through a single deterministic chokepoint (the ExecutionPlanFactory), and verifies the resulting answer through deterministic guardrails followed by a language-model judge. The architectural invariants on which the discipline rests are enforced by anti-drift tests rather than by review discipline alone --- a property that distinguishes the platform from a prototype.

\textit{The implementation across four cooperating microservices.} The manager service delivers the lifecycle backbone with its three-gate stage progression; the copilot service ships its first vertical slice through the v2 lane (Paths~1, 4, and the Path~8 safety family active end-to-end); the intelligence service delivers deterministic MR-package and workbook parsing with anomaly surfacing; the user interface integrates the backend services with role-aware presentation and dual portfolio / RFQ-bound copilot entry points.

\textit{The validation evidence.} Pytest suites across the backend services (259 / 636 / 64+48fd+6), anti-drift guards on the copilot's structural invariants, pipeline tests on its runtime invariants, and a structured business-impact mapping against the twelve audit pain points (2 addressed, 7 partial, 3 not yet) together support the claim that the architecture behaves as designed.

The limitations are stated honestly in Chapter~6. They cluster into \textit{validation gaps} (CI scoped to the manager service, SQLite rather than PostgreSQL execution, fixture-dependent parsers, no browser end-to-end tests, no curated language-model benchmark), \textit{architecture deferrals} (dedicated IAM service, working-memory injection into Planner/Compose/Judge, copilot-to-intelligence connector, copilot paths beyond 1, 4, and~8), and \textit{implementation hardening directions} (concurrency-safe RFQ code generation, cross-service correlation propagation, UI containerisation). None requires reopening of the architectural commitments developed in Chapter~4.

Future work follows directly. In the \textbf{near term}, the priority is hardening: CI across remaining services, PostgreSQL execution of the manager suite, recorded Newman runs, browser-level end-to-end tests, frontend cutover to the v2 thread-management lane, and a curated golden-answer benchmark for live-model evaluation. In the \textbf{medium term}, copilot path coverage extends to Paths~2, 3, 5, 6, and~7, the autonomous manager-to-intelligence cascade replaces manual triggers, the intelligence connector is wired, and working-memory injection is delivered under the trust-boundary discipline. In the \textbf{longer term}, the dedicated IAM service is extracted, and the Pillar~4 predictive layer is built on the historical dataset the lifecycle backbone is now beginning to accumulate.

This project leaves behind a platform whose foundations are sound, whose central architectural innovation is verifiable, and whose direction is established. It does not claim to have solved the problem identified at \clientCompany; it claims to have built the connective tissue on which that problem can be addressed over a multi-year horizon, under the engineering discipline that an industrial estimation context warrants.
